% Generated draft: requires mathematical and source review.
\documentclass[11pt]{article}
\usepackage[margin=25mm]{geometry}
\usepackage{fontspec}
\setmainfont{Latin Modern Roman}
\usepackage{amsmath,amssymb,mathtools,unicode-math}
\setmathfont{Latin Modern Math}
\usepackage{enumitem,needspace,xurl,hyperref}
\setlength{\parindent}{0pt}
\setlength{\parskip}{5pt}
\setlength{\emergencystretch}{3em}
\newcommand{\N}{\mathbb{N}}
\newcommand{\Z}{\mathbb{Z}}
\newcommand{\Q}{\mathbb{Q}}
\newcommand{\R}{\mathbb{R}}
\newcommand{\C}{\mathbb{C}}
\newcommand{\F}{\mathbb{F}}
\newcommand{\A}{\mathbb{A}}
\newcommand{\PP}{\mathbb{P}}
\newcommand{\E}{\mathbb{E}}
\newcommand{\eps}{\varepsilon}
\newcommand{\dd}{\,\mathrm{d}}
\begin{document}
\section*{UnsolvedMath ID: 2632}
\textbf{Source problem code:} catalogue entry
\subsection{Definitions and mathematical statement}
Under the usual left-pre-Lie convention, a pre-Lie ring is an abelian group \((R,+)\) with a biadditive operation \(\cdot\) satisfying
\[
(a\cdot b)\cdot c-a\cdot(b\cdot c)
=(b\cdot a)\cdot c-b\cdot(a\cdot c)
\qquad(a,b,c\in R).
\]
The convention is explicit because the supplied record does not reproduce Definition 66 or state its handedness.

The question attributed to A. Smoktunowicz is whether the specific extension of the Agrachev--Gamkrelidze construction proposed in Definition 66 of the cited 2025 paper produces a group from every pre-Lie ring in its intended domain. This is a question about that particular construction, not the assertion that an arbitrary operation on a pre-Lie ring is associative.

The supplied record provides no formula for the constructed multiplication, no explicit underlying constructed set, and no hypotheses from Definition 66. Consequently the construction cannot be replaced here by a guessed exponential, Baker--Campbell--Hausdorff formula or brace multiplication. In particular, possible nilpotence, torsion, coefficient, divisibility or convergence conditions are not supplied and must be recovered from that definition before one can state all input quantifiers precisely.

Once an underlying set \(X\) and proposed operation \(\circ\) from Definition 66 are actually fixed, the requested conclusion is that they satisfy
\[
\forall x,y,z\in X,\qquad (x\circ y)\circ z=x\circ(y\circ z),
\]
\[
\exists e\in X\ \forall x\in X:\quad e\circ x=x\circ e=x,
\qquad
\forall x\in X\ \exists y\in X:\quad x\circ y=y\circ x=e.
\]
These are only the target group axioms, not an invented definition of the missing operation.

The source further states conditionally that if this construction gives a group, it also gives a brace and hence applications to set-theoretic Yang--Baxter solutions. A left brace, in the common convention, has an abelian additive group and a multiplicative group satisfying \(a\circ(b+c)+a=a\circ b+a\circ c\); the record does not reproduce the proposed compatibility proof, so its brace consequence remains an attributed conditional statement. The supplied triage dated 2026-08-17 labels the question open, reports no verified proof that the extension always yields a group, and asks for associativity/inverse checks. The registry also labels it open. Conditional possible applications are not counted as a verified partial construction theorem.
\subsection{Short English statement}
Does the precise pre-Lie-to-group extension in Smoktunowicz's Definition 66 always satisfy the group axioms? The supplied record points to that definition but does not give its operation or hypotheses, so those cannot be invented. Its brace and Yang--Baxter consequences remain conditional on validity of that exact construction.
\subsection{Sources}
\begin{itemize}
\item Local UnsolvedMath registry, permanent ID 2632, source code catalogue entry, ``Kourovka Notebook Problem 21.123,'' from \emph{Kourovka Notebook}, New Problems, 21st issue (2026), original problem 21.123; attributed to A. Smoktunowicz. Includes supplied literature triage dated 2026-08-17.
\item A. Agrachev and R. Gamkrelidze, \emph{Journal of Soviet Mathematics} 17 (1981), 1650--1675, supplied reference for the original construction; no paper title, DOI or URL is supplied.
\item A. Smoktunowicz, \emph{Journal of Pure and Applied Algebra} 229, no. 12 (2025), article 108128, Definition 66; no paper title, DOI, URL or text of the definition is supplied.
\item \emph{Kourovka Notebook}, 21st issue (2026), Problem 21.123, supplied maintained source: \url{https://www.kourovka-notebook.org/}.
\end{itemize}
\end{document}
